\scanpage{97}
\fbox{Critères d'équivalence}

$X$ variété projective \deleted{\unsure{géométriquement} \unsure{normale}} \inserted{non sing} de dim $n$

$Y$ section hyperplane \deleted{\unsure{normale}} non sing

$X_m$ m.ième voisinage infinitésimal de $X$

\diagram{\left\{\begin{array}{ll}
\mathrm{Pic}(X)\longrightarrow\mathrm{Pic}(Y)&
\begin{array}{l}
\overbrace{\text{inj si }n\geq3,\ \mathrm{H}^1(Y,\mathcal{O}_Y(\unsure{-m}))=0\text{ pour }m>0;\ \text{surjectif si }n\geq4,\ \mathrm{H}^2(Y,\mathcal{O}_Y(\unsure{-m}))=0\text{ pour }m>0}^{\text{par \unsure{Lef-Grauert}}}\\
\uline{\text{noyau un }p\text{-groupe fini si }n\geq3\ (\text{nul si }Y\text{ général})};\quad\boxed{\text{noyau nul si }n\geq4}
\end{array}\\[1em]
\mathrm{Pic}^{\tau}(X)\longrightarrow\mathrm{Pic}^{\tau}(Y)&
\underbrace{\text{inj. exc. }p\text{ si }n\geq2,\ \text{bij. exc }p\text{ si }n\geq3;\ \text{injectif si }n\geq2\text{ et }Y\ \uline{\text{général}}}_{\text{via }\pi_1,\ \text{\unsure{Lef. Grauert}}}\\[1em]
\mathrm{Pic}^{0}(X)\longrightarrow\mathrm{Pic}^{0}(Y)&
\overbrace{\text{inj exc. }p\text{ si }n\geq2,\ \text{surjectif si }n\geq3,\ \text{bijectif si }n\geq2\text{ et }Y\ \uline{\text{général}}}^{\text{via }\unsure{\pi_1|\pi_2}\quad\text{\unsure{Lef. Grauert}}}
\end{array}\right.}
\margin{$\longleftarrow$ Pic des schémas non \unsure{réduits} \unsure{tels} $X_m$}

\diagram{\begin{array}{lll}
\mathrm{TN}(X)\longrightarrow\mathrm{TN}(Y)&\text{bij exc. }p\text{ si }n\geq3\quad\left[\begin{array}{l}\text{Si }n\geq3,\text{ le noyau est}\\\text{l'image de celui de}\\\mathrm{Pic}^{\tau}(X)\longrightarrow\mathrm{Pic}^{\tau}(Y)\end{array}\right]&\begin{array}{l}\uline{\text{injectif}}\text{ si }Y\ \uline{\text{général}}\\\text{et }n\geq3,\text{ ou}\\\text{si }n\geq4\end{array}\\[1em]
\mathrm{LN}(X)\longrightarrow\mathrm{LN}(Y)&\text{injectif si }n\geq3&\\[1em]
\mathrm{N}(X)\longrightarrow\mathrm{N}(Y)&\begin{array}{l}\text{noyau un }p\text{-groupe fini si }n\geq3\quad\left[\begin{array}{l}\text{contenu exc. }p\\\text{dans l'image}\\\text{du noyau }\deleted{\text{de}}\\\mathrm{Pic}(X)\longrightarrow\mathrm{Pic}(Y),\text{ qui est}\\\text{un }p\text{-groupe fini}\end{array}\right]\end{array}&\begin{array}{l}\text{nul si on a}\\n\geq4\\\uline{\text{ou}}\ Y\ \uline{\text{général}}\end{array}\\[1em]
\underset{m\text{ grand}}{\mathrm{Pic}(X)\longrightarrow\mathrm{Pic}(X_m)}&\text{inj si }n\geq3,\text{ bijectif si }n\geq4&
\end{array}}

\uline{Question} Surjectivité \inserted{(donc bijectivité)} de $\mathrm{Pic}(X)\longrightarrow\mathrm{Pic}(Y)$ \deleted{\unsure{pour}} pour $n\geq4$ sans \uline{conditions cohomologiques}, [\uline{équivalent} à celle de $\deleted{\unsure{T}}\mathrm{N}(X)\longrightarrow\deleted{\unsure{T}}\mathrm{N}(Y)$, ou encore celles de $\mathrm{TN}(X)\longrightarrow\mathrm{TN}(Y)$ et $\mathrm{LN}(X)\longrightarrow\mathrm{LN}(Y)$]. Équivalent à celle de $\mathrm{Pic}(X_m)\longrightarrow\mathrm{Pic}(Y)$ pour $m$ grand. N.B. la surjectivité mod $p^m$ (car. $p>0$) \inserted{[regarder les obstructions successives, \unsure{s'annulent} par $p$\ldots]} \uline{donne pour} $n\geq4$, $\mathrm{Pic}(X)\longrightarrow\mathrm{Pic}(Y)$ \uline{est surjectif exc.} $p$. \deleted{\unsure{Que}} En car. 0, \uline{on peut invoquer la théorie Kählérienne}, qui prouve que \deleted{\illegible{2}} $\mathrm{N}\deleted{\unsure{^{\tau}}}(X)\xrightarrow{\sim}\mathrm{N}(Y)$, d'où $\mathrm{Pic}(X)\longrightarrow\mathrm{Pic}(Y)$.

\margin{Corollaire : $\rho(X)=\rho(Y)$ pour $n\geq4$ [$\rho(X)\leq\rho(Y)$ si $n\geq3$]}
\margin{\deleted{$\unsure{\rho}(X)\leq\rho(Y)$ si $n\geq2$ et $Y$ général}}

\uline{Démonstration} algébrique ???
\margin{96}
